\chapter{Les nombres décimaux}\label{ch:g5:decimals}

Les dixièmes et les centièmes sont apparus dans
\cref{ch:g4:decimals}~; un niveau de plus, les millièmes, complète
le système décimal. Ce chapitre consolide la lecture, la
décomposition et --- surtout --- la comparaison des nombres
décimaux sans tomber dans leurs pièges célèbres.

\section{Jusqu'aux millièmes}

\begin{definition}[Positions décimales]\label{def:g5:decimals:places}
Après la \omterm{def:g4:decimals:point}{virgule} viennent les dixièmes, les centièmes, les
millièmes --- chacun valant dix fois moins que le précédent~:
\[
4.362 = 4 + \frac{3}{10} + \frac{6}{100} + \frac{2}{1000}
= \frac{4\,362}{1000}.
\]
\end{definition}

\begin{omfigure}
\begin{tikzpicture}[scale=0.9]
  \draw (0,0) grid[xstep=2.3, ystep=0.85] (9.2,1.7);
  \node[font=\small] at (1.15,1.275) {unités};
  \node[font=\small] at (3.45,1.275) {dixièmes};
  \node[font=\small] at (5.75,1.275) {centièmes};
  \node[font=\small] at (8.05,1.275) {millièmes};
  \node[font=\large, omDef] at (1.15,0.425) {$4$};
  \node[font=\large, omThm] at (3.45,0.425) {$3$};
  \node[font=\large, omThm] at (5.75,0.425) {$6$};
  \node[font=\large, omThm] at (8.05,0.425) {$2$};
  \fill (2.3,0.06) circle (2.5pt);
\end{tikzpicture}

{\small Le tableau de numération de $4.362$~: partie entière en
bleu, partie décimale en rouge, la \omterm{def:g4:decimals:point}{virgule} entre unités et
dixièmes.}
\end{omfigure}

\begin{example}[Plusieurs lectures, un seul nombre]\label{ex:g5:decimals:readings}
$4.362$ peut se lire «~quatre \omterm{def:g4:decimals:point}{virgule} trois six deux~», ou «~quatre
et $362$ millièmes~», ou se décomposer en $4 + 0.3 + 0.06 + 0.002$.
Et ajouter des \omterm{def:g1:counting:zero}{zéros} finaux ne change rien~: $4.362 = 4.3620$.
Mais $4.362 \neq 4.0362$~: un \omterm{def:g1:counting:zero}{zéro} \emph{à l'intérieur} pousse
chaque chiffre vers une place plus petite.
\end{example}

\begin{example}[Où est-il sur la droite~?]\label{ex:g5:decimals:line}
Pour placer $4.362$~: entre $4$ et $5$~; en zoomant, entre $4.3$ et
$4.4$~; en zoomant encore, entre $4.36$ et $4.37$, plus près de
$4.36$. Chaque chiffre décimal est un niveau de zoom de plus.
\end{example}

\begin{omfigure}
\begin{tikzpicture}[scale=1.15]
  \draw[-Stealth] (-0.2,1.6) -- (10.4,1.6);
  \foreach \i in {0,...,10} \draw (\i,1.52) -- (\i,1.68);
  \node[below, font=\small] at (0,1.5) {$4$};
  \node[below, font=\small] at (10,1.5) {$5$};
  \node[below, font=\small] at (3,1.5) {$4.3$};
  \node[below, font=\small] at (4,1.5) {$4.4$};
  \fill[omThm] (3.62,1.6) circle (1.7pt);
  \draw[omExo, thin] (3,1.45) -- (0,0.15);
  \draw[omExo, thin] (4,1.45) -- (10,0.15);
  \draw[-Stealth] (-0.2,0) -- (10.4,0);
  \foreach \i in {0,...,10} \draw (\i,-0.08) -- (\i,0.08);
  \node[below, font=\small] at (0,-0.1) {$4.3$};
  \node[below, font=\small] at (10,-0.1) {$4.4$};
  \node[below, font=\small] at (6,-0.1) {$4.36$};
  \fill[omThm] (6.2,0) circle (1.7pt)
    node[above, font=\small] {$4.362$};
\end{tikzpicture}

{\small Zoom sur la droite graduée~: le segment de $4.3$ à $4.4$,
agrandi, est lui-même découpé en dix centièmes --- et $4.362$ se
situe juste après $4.36$.}
\end{omfigure}

\section{Comparer des décimaux}

\begin{method}[Comparer deux décimaux]\label{met:g5:decimals:compare}
\begin{enumerate}
  \item Comparer les parties entières~: $6.1 > 5.987$.
  \item Parties entières égales~: comparer les parties décimales
        chiffre par chiffre depuis la \omterm{def:g4:decimals:point}{virgule} --- dixièmes, puis
        centièmes, puis millièmes~;
  \item compléter avec des \omterm{def:g1:counting:zero}{zéros} finaux rend la comparaison
        équitable~: $2.5$ contre $2.48$ devient $2.50$ contre
        $2.48$, et $50 > 48$ centièmes.
\end{enumerate}
Le piège, une dernière fois~: \emph{une partie décimale plus longue
ne signifie pas un nombre plus grand} --- $2.48$ a plus de chiffres
que $2.5$ et est plus petit.
\end{method}

\begin{example}\label{ex:g5:decimals:order}
Ranger $7.3$~; $7.09$~; $7.31$~; $7.299$. Compléter à trois
décimales~: $7.300$~; $7.090$~; $7.310$~; $7.299$. Alors
\[
7.09 < 7.299 < 7.3 < 7.31 .
\]
Remarquer $7.299 < 7.3$ même si $299$ semble grand~: $299$
millièmes contre $300$ millièmes.
\end{example}

\begin{example}[Se glisser entre deux décimaux]\label{ex:g5:decimals:between}
Y a-t-il un nombre entre $5.7$ et $5.8$~? Oui, plein~: $5.75$,
$5.71$, $5.799$\,\dots\ Entre $5.79$ et $5.8$~? Encore plein~:
$5.795$, par exemple. Entre deux décimaux différents, on peut
\emph{toujours} en glisser un autre --- il n'y a pas de décimal
«~suivant~».
\end{example}

\section{Arrondir des décimaux}

\begin{definition}[Arrondi]\label{def:g5:decimals:round}
Pour \omterm{def:g4:numbers:round}{arrondir} à l'unité (ou au dixième, ou au centième), on regarde
le chiffre suivant~: $5$ ou plus arrondit vers le haut, sinon vers
le bas. Ainsi $4.362$ s'arrondit à $4$ (unité), à $4.4$
(dixième), à $4.36$ (centième).
\end{definition}

\begin{example}[L'argent s'arrondit aux centimes]\label{ex:g5:decimals:money}
Un prix calculé à $7.4983$ s'affiche $7.50$~: les montants de la
vie réelle s'arrondissent au centième. Noter la cascade~:
$7.4983 \to 7.50$, où le $49$ est devenu $50$ --- l'\omterm{def:g4:numbers:round}{arrondi} peut
se propager sur plusieurs chiffres.
\end{example}

\section{Exercices}

\begin{exercise}[$\star$]\label{exo:g5:decimals:1}
Écrire comme un nombre décimal~: $5 + \frac{2}{10} + \frac{7}{1000}$~;
\quad $\frac{85}{100}$~; \quad $\frac{4\,507}{1000}$~; \quad «~douze
et neuf millièmes~».
\end{exercise}

\begin{exercise}[$\star$]\label{exo:g5:decimals:2}
Dans $27.418$~: que compte le $4$~? Le $8$~? Décomposer le nombre
comme dans \cref{def:g5:decimals:places}.
\end{exercise}

\begin{exercise}[$\star$]\label{exo:g5:decimals:3}
Vrai ou faux~? Expliquer avec les positions.
$3.60 = 3.6$~; \quad $0.5 = 0.05$~; \quad $2.070 = 2.07$~; \quad
$1.3 = 1.03$.
\end{exercise}

\begin{exercise}[$\star$]\label{exo:g5:decimals:4}
Recopier et compléter avec $<$, $>$ ou $=$~:
\[
4.8 \;?\; 4.79, \qquad
0.302 \;?\; 0.32, \qquad
6.5 \;?\; 6.500, \qquad
9.09 \;?\; 9.1 .
\]
\end{exercise}

\begin{exercise}[$\star$]\label{exo:g5:decimals:5}
Ranger du plus petit au plus grand~: $3.14$~; \ $3.4$~; \ $3.041$~; \
$3.104$~; \ $3.41$.
\end{exercise}

\begin{exercise}[$\star$]\label{exo:g5:decimals:6}
Entre quels deux entiers se situe chaque nombre~? Entre quels deux
nombres à une décimale~? \quad $6.83$~; \quad $0.492$~; \quad
$12.06$.
\end{exercise}

\begin{exercise}[$\star$]\label{exo:g5:decimals:7}
Arrondir $8.276$~: à l'unité~; au dixième~; au centième. Arrondir
$3.95$ au dixième.
\end{exercise}

\begin{exercise}[$\star$]\label{exo:g5:decimals:8}
Donner un nombre strictement entre~: $2.6$ et $2.7$~; \quad $0.41$
et $0.42$~; \quad $5.99$ et $6$.
\end{exercise}

\begin{exercise}[$\star$]\label{exo:g5:decimals:9}
Les hauteurs de quatre plantes sont $0.85$~m, $1.2$~m, $0.9$~m et
$1.02$~m. Les ranger de la plus courte à la plus haute.
\end{exercise}

\begin{exercise}[$\star\star$]\label{exo:g5:decimals:10}
Lucie dit~: «~$7.12 > 7.9$ parce que $12 > 9$~». Corriger son
erreur, en comparant d'abord les dixièmes.
\end{exercise}

\begin{exercise}[$\star\star$]\label{exo:g5:decimals:11}
Trouver tous les nombres à exactement deux chiffres après la
\omterm{def:g4:decimals:point}{virgule} qui s'arrondissent à $6.4$ lorsqu'on arrondit au dixième.
Quels sont le plus petit et le plus grand~?
\end{exercise}
